% Build in this directory: pdflatex -interaction=nonstopmode -halt-on-error class8_handout.tex
\documentclass[a4paper,10pt,pdflatex,ja=standard,jafont=ipaex-type1,
  magstyle=real,scale=1]{bxjsarticle}
\usepackage{lmodern}
\usepackage{amsmath,amssymb,mathtools}
\usepackage{geometry}
\geometry{reset,left=22mm,right=22mm,top=20mm,bottom=17mm,
  headheight=13pt,headsep=6mm,footskip=10mm}
\usepackage{xcolor}
\usepackage{enumitem}
\usepackage{array,booktabs,tabularx}
\usepackage{fancyhdr}

\definecolor{ink}{HTML}{193A5C}
\definecolor{accent}{HTML}{237D82}
\definecolor{soft}{HTML}{EAF3F3}
\definecolor{muted}{HTML}{52616B}
\newcommand{\coursename}{解析学 II}
\newcommand{\handoutnumber}{8}
\pagestyle{fancy}
\fancyhf{}
\fancyhead[L]{\small\color{ink}\coursename{}・授業要約と演習}
\fancyhead[R]{\small\color{ink}第\handoutnumber 回}
\fancyfoot[C]{\small\color{ink}\thepage/2}
\renewcommand{\headrulewidth}{0.3pt}
\setlength{\parindent}{0pt}
\setlength{\parskip}{2pt}
\setlength{\abovedisplayskip}{5pt}
\setlength{\belowdisplayskip}{5pt}
\setlength{\abovedisplayshortskip}{3pt}
\setlength{\belowdisplayshortskip}{4pt}
\setlist[itemize]{leftmargin=1.8em,itemsep=1pt,topsep=2pt,parsep=0pt}

\newcommand{\handouttitle}[1]{%
  \begin{center}
    {\Large\color{ink}#1}\par\vspace{3mm}
    {\color{accent}\rule{0.88\linewidth}{0.7pt}}
  \end{center}\vspace{-2mm}}
\newcommand{\handout}[2]{%
  \renewcommand{\handoutnumber}{#1}%
  \handouttitle{第#1回\quad #2}}
\newcommand{\topic}[1]{\par\vspace{5pt}%
  {\large\sffamily\mdseries\color{accent}#1}\par\vspace{2pt}}
\newcommand{\keybox}[1]{\par\vspace{3pt}\begingroup
  \setlength{\fboxsep}{4pt}%
  \colorbox{soft}{\parbox{\dimexpr\linewidth-2\fboxsep\relax}{\small #1}}%
  \endgroup\par\vspace{2pt}}
\newenvironment{problems}
  {\begin{enumerate}[leftmargin=2em,itemsep=3pt,topsep=3pt,parsep=0pt,
    font=\color{accent}]}
  {\end{enumerate}}


\newcommand{\examplelabel}[1]{{\sffamily\mdseries\color{ink}#1}}
\newcommand{\answerroom}{\par\vspace{16mm}}

\begin{document}
\fontsize{10pt}{13.5pt}\selectfont
\setlength{\abovedisplayskip}{4pt}
\setlength{\belowdisplayskip}{4pt}
\handout{8}{前半まとめ：第1回から第7回}
\topic{逆関数と逆三角関数}
一対一対応 $f:A\to B$ の逆関数を $g:B\to A$ とすると
$$
 f(g(y))=y,\quad g(f(x))=x,\qquad
 g'(y)=\frac1{f'(g(y))}.
$$
微分公式では $A,B$ を開区間、$f$ を $C^1$ 級、$f'(g(y))\ne0$ とする。
逆三角関数の値域は $\arcsin x\in[-\pi/2,\pi/2]$、$\arccos x\in[0,\pi]$、
$\arctan x\in(-\pi/2,\pi/2)$ である。
$$
 (\arcsin x)'=\frac1{\sqrt{1-x^2}},\quad
 (\arccos x)'=-\frac1{\sqrt{1-x^2}}\quad(|x|<1),\qquad
 (\arctan x)'=\frac1{1+x^2}.
$$

\topic{有理関数の積分}
多項式除法 → 分母の実因子分解 → 部分分数分解 → 基本公式で積分する。
一次因子の項には定数分子、既約二次因子の項には一次式の分子を置く。
$$
 \int\frac{dx}{x-a}=\log|x-a|+C,\qquad
 \int\frac{dx}{x^2+c^2}=\frac1c\arctan\frac xc+C\quad(c>0).
$$

\topic{重積分の定義と累次積分}
重積分は $\sum_Qf(\xi_Q)\operatorname{Area}(Q)$ の、代表点によらない細分極限である。
長方形上の連続関数は可積分。一般の有界領域では境界の条件も確認する。
$$
 D=\{(x,y):a\le x\le b,\ \phi(x)\le y\le\psi(x)\}
 \ \Longrightarrow\ 
 \iint_Df=\int_a^b\left(\int_{\phi(x)}^{\psi(x)}f(x,y)\,dy\right)dx.
$$
同じ領域を横に切る形へ書き直せば、積分順序を交換できる。

\topic{広義積分：特異点と無限遠}
特異点の各側で極限を別々に取る。符号が変わる関数では絶対収束を確認する。
$$
 \int_0^1x^{-p}\,dx<\infty\iff p<1,\qquad
 \int_1^\infty x^{-p}\,dx<\infty\iff p>1.
$$
非負関数の積分は、領域を単調に増やした極限で扱える。

\topic{変数変換とヤコビアン}
$D=T(E)$ の両領域は有界で境界の面積が0とする。
$T(u,v)=(x(u,v),y(u,v))$ が $E$ の近傍で $C^1$ 級、内部で1対1かつ $\det DT\ne0$ なら
$$
 \iint_{T(E)}f(x,y)\,dx\,dy
 =\iint_E f(T(u,v))\,|\det DT(u,v)|\,du\,dv.
$$
極座標では $x=r\cos\theta,\ y=r\sin\theta$、$dx\,dy=r\,dr\,d\theta$。
円・扇形・円環では半径と角度の範囲を先に決める。

\topic{体積・密度・総量}
$V=\{(x,y,z):(x,y)\in D,\ \phi\le z\le\psi\}$、$\phi\le\psi$ のとき
$$
 \operatorname{Vol}(V)=\iint_D(\psi-\phi)\,dx\,dy,\qquad
 M=\iiint_V\rho\,dV.
$$
体積・質量の計算でも、積分領域と座標変換を正しく書くことが出発点となる。
\keybox{公式を選ぶ前に、定義域・特異点・領域の形・可積分性を点検する。計算結果だけでなく設定の理由も説明する。}

\newpage
\handouttitle{第8回\quad 演習問題}
計算過程を記し、必要な条件・積分領域・向きを明示すること。
\begin{problems}
\item $f(x)=x^2+2x$（$x\ge-1$）の逆関数とその定義域を求めよ。また $\arcsin(-1/2)$ と $\arccos(-1/2)$ を求めよ。\answerroom
\item $\displaystyle\int\frac{x^2+1}{x(x+1)}\,dx$ を求めよ。多項式部分と部分分数の項を明記すること。\answerroom
\item $[0,1]^2$ 上の連続関数 $f$ の重積分をリーマン和で定義せよ。$f(x,y)=xy$ について、累次積分で値を計算せよ。\answerroom
\item $D=\{(x,y):0\le y\le x\le1\}$ 上で $\iint_D(x^2+y^2)\,dx\,dy$ を求めよ。順序を交換した式も書け。\answerroom
\item $p\in\mathbb R$ に対して $\displaystyle\int_0^\infty\frac{dx}{(1+x)^p}$ の収束条件と、収束する場合の値を求めよ。\answerroom
\item $a,b>0$、$D=\{(x,y):x^2/a^2+y^2/b^2\le1\}$ とする。$\iint_D(x^2/a^2+y^2/b^2)\,dx\,dy$ を変数変換で求めよ。\answerroom
\item $V=\{(x,y,z):x^2+y^2\le1,\ 0\le z\le2\}$、$\rho=1+z$ の質量を求めよ。円柱座標で積分を設定すること。\answerroom
\end{problems}
\end{document}
